\section{Методика и организация экспериментов}
\label{sec:methodology}

\subsection{Цель экспериментальной части}

Экспериментальная часть должна ответить на вопрос, подтверждается ли
выдвинутая гипотеза -- в части возможности потокового обновления
in-memory базы после выделения микросервиса. Для этого нужно измерить,
как новое и прежнее внутренние хранилища ведут себя на наборе операций, соответствующих
реальным сценариям использования сервиса:

\begin{itemize}
\item полная выгрузка таблицы (\texttt{Dump});
\item одиночное обновление маршрута (\texttt{Upsert});
\item точечный поиск по IP-адресу (\texttt{Lookup});
\item влияние числа шардов на ключевые операции;
\item различия между IPv4 и IPv6.
\end{itemize}

\subsection{Сравниваемые варианты}

В опытах сравниваются:

\begin{itemize}
\item \textbf{\texttt{ranger}} -- прежний вариант хранилища на базе
библиотеки \texttt{ranger};
\item \textbf{шард. ART} -- новый вариант хранилища на базе шардированного
Adaptive Radix Tree.
\end{itemize}

Это важный методический момент: сравнение выполняется не между двумя
абстрактными структурами данных из учебника, а между двумя реальными
вариантами внутреннего хранилища одного и того же сервиса с одним и тем же
\texttt{prefixStore}-контрактом. Все различия, фиксируемые в
эксперименте, можно интерпретировать как различия между
реализацией, ориентированной на полную замену, и реализацией,
ориентированной на поток обновлений, в рамках этого
контракта.

\subsection{Конфигурация стенда}

Параметры экспериментального стенда приведены в
табл.~\ref{tab:bench-setup}. Все измерения проводились на одной и той
же физической машине, без использования контейнеров и виртуализации,
с зафиксированным набором фоновых процессов.

\begin{table}[ht]
\centering
\caption{Конфигурация стенда экспериментов}
\label{tab:bench-setup}
\small
\begin{tabular}{>{\raggedright\arraybackslash}p{4.5cm}>{\raggedright\arraybackslash}p{8.5cm}}
\toprule
Параметр & Значение \\
\midrule
CPU                  & AMD Ryzen AI 9 HX 370 (24 потока) \\
ОЗУ                  & 32 ГБ LPDDR5X \\
Среда                & Linux x86\_64 \\
Среда исполнения     & Go 1.24+ \\
Инструмент           & \texttt{go test -bench} с теговой подсистемой \texttt{bench}, \texttt{benchheavy} \\
ART-параметры        & 16 шардов, \texttt{WithDumpSort(false)} в сценариях с большими $n$ \\
	Размеры таблиц       & $n \in \{1\,000;\, 10\,000;\, 100\,000\}$; отдельный тяжёлый сценарий -- $3\,000\,000$ смешанных IPv4/IPv6 записей \\
	Семейства адресов    & IPv4 \texttt{/24} (искусственные диапазоны \texttt{10.x.y.0/24}, \texttt{11.x.y.0/24}), IPv6 \texttt{/64} в \texttt{2001:db8::/32}; для тяжёлого сценария -- смесь IPv4 \texttt{/24} и IPv6 \texttt{/64} \\
Параметры запуска    & \texttt{count=2}, \texttt{benchtime=400ms}\,/\,\texttt{500ms} \\
\bottomrule
\end{tabular}
\end{table}

\subsection{Наборы данных и сценарии}

Для измерений используются синтетические наборы префиксов, а также
сценарии, имитирующие реалистичную нагрузку. В частности:

\begin{itemize}
\item таблицы различного размера, вплоть до~$10^5$ префиксов в основных
сериях и отдельный тяжёлый сценарий на $3\cdot10^6$ смешанных записей;
\item отдельные наборы для IPv4 и IPv6;
\item сценарии полного \texttt{Dump};
\item сценарии одиночного \texttt{Upsert} поверх уже загруженной
таблицы;
\item сценарии \texttt{Lookup};
\item sweep по числу шардов $K \in \{1,8,16,32,64\}$.
\end{itemize}

Такой дизайн эксперимента позволяет отдельно оценить операции, которые
ограничивают применимость каждой реализации: полную выгрузку, точечный
поиск, инкрементальное обновление и влияние числа шардов.

\subsection{Почему именно эти метрики достаточны}

Для данной работы важно подчеркнуть, что метрики выбираются не
случайно, а отражают конкретные эксплуатационные сценарии.

Операция \texttt{Dump} соответствует сценарию получения полного снимка
внешним потребителем, а также внутренним этапам начальной инициализации и
ресинхронизации.

Операция \texttt{Upsert} соответствует основной сущности потокового
режима: локальному отражению одного BGP-изменения. Именно её стоимость
определяет, осмыслен ли переход от полной пересборки таблицы к потоку;
измерение \texttt{Upsert} -- прямая проверка этой части гипотезы, а не
«бенчмарк ради ART».

Операция \texttt{Lookup} задаёт дополнительный разрез: приемлемость
задержки точечного запроса при выбранном хранилище (включая реализацию
LPM).

Измерение при разных значениях числа шардов показывает, где проходит
граница между полезной параллельностью и издержками чрезмерного дробления
состояния.

Именно эта совокупность метрик позволяет проверить, достигнута ли
цель перехода к потоковой архитектуре, и подтвердить или
опровергнуть выдвинутую гипотезу.

\subsection{Ограничения экспериментальной методики}

Экспериментальная часть измеряет прежде всего внутренние операции
хранилища: полную загрузку, точечный поиск, инкрементальное обновление и
полную выгрузку таблицы. Такой выбор позволяет изолировать различия между
двумя реализациями \texttt{prefixStore}, но оставляет за пределами
измерений ряд факторов эксплуатации:

\begin{itemize}
\item сетевые задержки в внешних gRPC-вызовах;
\item долгосрочную динамику при реальном потоке BGP update-событий за
большие интервалы времени;
\item поведение под многоклиентской нагрузкой на пользовательский API;
\item влияние сборщика мусора и давления на аллокатор при длительной
работе сервиса;
\item воспроизведение реалистичного распределения длин префиксов из
промышленных BGP-таблиц.
\end{itemize}

Для цели данной работы такое ограничение методики приемлемо: измерения
проверяют именно тот слой, в котором различаются \texttt{ranger} и
шардированный ART, и позволяют сравнить два варианта хранилища по
ключевым операциям.
Ограничения интерпретации результатов перечислены выше и учитываются в
выводах главы~\ref{sec:experiment-results}.

%==============================================================
